\chapter{Development Process Analysis} % {{{
\label{cha:Development Process Analysis}

\section{Project Analyses Criteria} % {{{
\label{sec:Project Analyses Criteria}

% }}}

\cleardoublepage
\input{chapters/03_analysis_drupal}

\cleardoublepage
\input{chapters/03_analysis_plone}

\cleardoublepage
\input{chapters/03_analysis_python}

\cleardoublepage
\input{chapters/03_analysis_php}

\cleardoublepage
\input{chapters/03_analysis_gnome}

\cleardoublepage
\input{chapters/03_analysis_kde}

\cleardoublepage
\section{Definition and Origin of the Catalogue} % {{{
\label{sec:Definition and Origin of the Catalogue}

To characterize the development process of FOSS project, the following
catalogue was established. By covering all listed points one a project should
have been analysed thoroughly and lead to a reasonable breakdown of all
projects.

\begin{description}

  \item[Description of the Project] A general description about the project
    should be provided to inform the reader about the goals and current state
    of a project.

  \item[Project Category] In order to provide methods to compare projects, the
    project category can help to contrast projects of the same and other
    categories. Examples are programming languages or content management
    systems.

  \item[Scope of Analysis] As FOSS projects often are hard to analyze due to
    their size and unclear definition of modules, the scope will be limited to
    a well known subset of the whole project.

  \item[License] Some projects use their own license and some use the well
    known General Public License. It is interesting to find out whether this
    makes any difference in the development process of a module.

  \item[History] As only projects with a quite long history are analyzed it
    makes sense to glance at the project's history.

  \begin{description}

    \item[Founders] In some projects, the original founders are still present
      and have a influential voice, in others the original authors are no
      longer involved. To analyzed this further, the authors have to be
      introduced and their role will be analyzed in the community part.

    \item[Project Age] Projects always need time to evolve and to find their
      best development process. As this needs time, the project age can give
      some information in what development process state this project is and
      how it will evolve in the future.

  \end{description}

  \item[Community] The people behind the project are the driving force. Without
    them, a project would not exist. Therefore it is important to analyze the
    diverse community of a project.

  \begin{description}

    \item[Size of the Community] An important measure is the approximate size
      of a community. Several structural changes inside the project can depend
      on this size.

    \item[Communication inside the Project] The communication is a vital thing
      in an open project. It is interesting to compare the different methods of
      communication in several projects depending on their development
      structure and size.

    \item[Conferences and Meet-ups] Meetings of developers are an important
      element of the development process in FOSS projects. Also it shows, that
      there is enough interest available to provide the money needed for
      preparing such events.

    \item[Roles] The development process is often defined and lead by important
      roles in the project. Also, depending on what role a person has in the
      project, his influence varies. 

    \item[Role of the Founders] In some projects the founders are still
      actively involved and in some not. Depending on that, the founders often
      have a very important role in the project development process.

  \end{description}

  \item[Release Process] In the FOSS world an often heard statement is that a
    new version will be released when it is ready. In fact however quite all
    projects have a certain release schedule or process they follow.

  \begin{description}

    \item[Version Naming] Each project provides a specific version naming
      scheme to which they adopt and which characterizes major, minor or bug
      fixing releases. This information is vital to understand the project's
      release schedule.

    \item[Characterization of Major and Minor Releases] Most projects provide
      some kind of major releases with new features which is often incompatible
      with previous releases, minor releases which mostly are backwards
      compatible and provide new features and finally bug fixing only releases.
      This however differs from project to project.

    \item[Release Schedule] Even if there are no fixed release cycles, most
      projects have nevertheless a release schedule in action. The release
      schedule defines deadlines and important steps in the process of creating
      a new release.

    \item[Important Steps in the Schedule] These often define freezes or
      specific points in time when all members of the project are restricted to
      for example not provide new code in order to increase the stability of
      the new release.

  \end{description}

  \item[Development] The actual development is closely interweaved with the
    release process and defines how and when the development of a project takes
    place.

  \begin{description}

    \item[Development Lead] In fact, most FOSS projects are lead by groups of
      people who define new features and lead the development process.
      Additionally there are hardly any projects which have a completely
      unstructured development process and have no development leaders in
      place.

    \item[Development Workflow] The actual development process as often defined
      by the project leaders provides ways and methods to propose and develop
      new features for the upcoming release.

    \item[Feature Inclusion Process] This inclusion process is used by some
      projects which defines process how and when new features and come into
      the project.

  \end{description}

\end{description}

% }}}

\section{PostgreSQL Project Analysis} % {{{

\marginpar{\includegraphics[width=\marginparwidth]{postgresql/logo}}

The PostgreSQL project provides an \ac{ORDBMS} which runs on all major systems
such as Linux, UNIX, Microsoft Windows or Mac OS X
\cite{PostgreSQLAbout,PostgreSQLFAQ}. The group behind the project is known as
PostgreSQL Global Development Group which consists of several volunteers and
employed developers. PostgreSQL is quite popular amongst it's users and has won
several prizes for the best database management system \cite{PostgreSQLAwards}.
According to the project it is the leading \ac{FOSS} database system with
thousands of users and contributors \cite{PostgreSQLPressKit}.

\begin{description}

  \item[Project Category] PostgreSQL is a database management system, more
    specifically a \ac{ORDBMS} \cite{PostgreSQLAbout}.

  \item[Scope of Analysis] The PostgreSQL Core Distribution will be analyzed,
    which consists of the PostgreSQL database server, several tools and
    bindings around it \cite{PostgreSQLDownload}.

  \item[License] The PostgreSQL project makes use of the PostgreSQL license,
    which is a \ac{FOSS} license \cite{PostgreSQLFAQ,PostgreSQLLicense}. It
    only requires to maintain the copyright and licensing information in the
    licensed source code and therefore it is quite similar to a BSD license.

  \item[History] The project was started in 1986 by Lawrence A. Rowe and
    Michael R. Stonebraker at the University of California in Berkeley under
    the name POSTGRES \cite{PostgreSQLHistory}. Not until 1996 it was developed
    in a university style fashion trying to explore new areas in database
    management systems. It was then released as \ac{FOSS} by two students of
    Stonebraker with the new name PostgreSQL the adopted \ac{SQL}. It then
    received a big development boost and emerged to one of the leading
    \acp{ORDBMS}.

  \item[Community] The PostgreSQL project has a quite large community available
    which will be described in the following.

  \begin{description}

    \item[Size of the Community] According to the project there are 6 core team
      members, 38 major contributors and 42 contributors
      \cite{PostgreSQLContributors}. Adding no longer active persons the number
      sums up to around 140. The number of total contributors might however be
      higher, as the project only includes people who have made contributions
      over a long time.

    \item[Communication inside the Project] The communication mostly happens
      over mailing lists, although also \ac{IRC} channels exist. The most
      important mailing list for development is the pgsql-hackers mailing list
      \cite{PostgreSQLDevFAQ}.

    \item[Conferences and Meet-ups] The most important international conference
      is the annual PgCon which first happened in 2007 \cite{PostgreSQLEvents}.
      However there are a lot of local conferences and workshops available.

    \item[Roles] The developers are split into three groups
      \cite{PostgreSQLContributors}. The core team decides on the general
      direction of the PostgreSQL project as well as the release cycle and the
      releases. Major contributors are people who introduced or maintain big
      features to the project. The contributors are all other who provide
      patches to the project.

    \item[Role of the Founders] Michael R. Stonebraker and Lawrence A. Rowe did
      not follow the project when it was published as \ac{FOSS} in 1996
      \cite{PostgreSQLHistory}. They are therefore no longer actively involved
      in the project.

  \end{description}

  \item[Release Process] The used release process is highly structured and will
    be described in the following.

  \begin{description}

    \item[Version Naming] The PostgreSQL project uses a three digit number
      scheme which consists of a major, minor and micro number
      \cite{PostgreSQLVersioning}.

    \item[Characterization of Major and Minor Releases] The incrementation of a
      major or minor number defines a major release including new features and
      often backwards incompatible changes \cite{PostgreSQLVersioning}. Such a
      release occurs roughly once a year
      \cite{PostgreSQLDevelopment,PostgreSQLFAQ}. For each major release there
      exist a number of minor releases which are defined by incrementing the
      micro version number. Only bug and security fixes are allowed for those
      releases. Each major release is supported for five years by the
      PostgreSQL project. However in some cases the project can decide to drop
      support for a specific release if bugs cannot be resolved without risking
      the stability.

    \item[Release Schedule] As major releases occur every year, the project
      uses the same release schedule every year along with some minor
      improvements from the last years \cite{PostgreSQLDevelopment}. The
      release schedule gets planned and decided during the PgCon conference. It
      basically starts each June with a commit review in which possible patches
      might be included in the next major release. Such reviews, also known as
      Commit Fest happen four times each two months apart. In the month between
      an alpha release is published. After the last Commit Fest, more alpha
      releases can follow, if not beta releases are published. If no more
      critical errors are found several release candidates will be published
      which will lead to the next major release.

    \item[Important Steps in the Schedule] The Commit Fest are the only
      possibility to add new features to the project and are therefore a vital
      part of the development process
      \cite{PostgreSQLDevelopment,PostgreSQLCommitFest}. 

  \end{description}

  \item[Development] The development of the PostgreSQL project is closely
    entangled with the release process and will be described in the following.

  \begin{description}

    \item[Development Lead] The PostgreSQL project is mostly driven by the core
      team \cite{PostgreSQLDevFAQ,PostgreSQLFAQ,PostgreSQLContributors}. Quite
      all decisions are made by them as well as most new features and code
      contributions.

    \item[Development Workflow] Small patches are often submitted directly to
      the repository if they are done by a major contributor or a core team
      member \cite{PostgreSQLDevFAQ}. All other patches however have to be
      reviewed and therefore be sent to the pgsql-hackers mailing list. New
      features are proposed in the Commit Fests.

    \item[Feature Inclusion Process] The already mentioned Commit Fest is a
      periodic break in which no new development is done, however already
      existing patches are reviewed and get feedback
      \cite{PostgreSQLCommitFest}. Not only the core team does those reviews
      but also other developers. A single Commit Fest usually runs for about
      one month explaining the gap of one month between each Commit Fest. One
      such event is lead by a Commit Fest Manager who is responsible that all
      submitted patches will be reviewed. Often, a review results in a
      discussion on the mailing list and follows an acceptance, a return with
      feedback or a rejection. A patch can have several states, depending on
      whether it is in progress or if it is finished
      \cite{PostgreSQLCommitFestRunning}. If a patch is in progress, the
      following states might apply.

      \begin{description}

        \item[Needs Review] A patch is not reviewed and waiting for a review.

        \item[Waiting on Author] A new version of the patch is expected by the
          author.

        \item[Discussing Review Results] The review was done, however it is
          discussed on the mailing list.

        \item[Ready for Committer] The review was done and no issues were
          found. It is now waiting for a final review.

      \end{description}

      The in progress states of a patch are the following.

      \begin{description}

        \item[Returned With Feedback] The patch was reviewed and feedback was
          given. A new version of the patch is expected for the next Commit Fest.

        \item[Rejected] The patch was rejected and won't make it into the next
          major release.

        \item[Committed] The patch was applied and no remaining issues are
          assured.

      \end{description}

  \end{description}

\end{description}

% }}}

% }}}
